\section{Methods}
\label{methods}

\subsection{Problem Formulation and Design Principle}

We consider closed-loop collaborative autonomous driving in CARLA, where an ego vehicle must follow a prescribed route while interacting with dynamic traffic participants, including vehicles, pedestrians, and bicycles. The upstream system (collaborative V2X perception and a learned CoDriving planner) provides waypoints and a desired speed at each decision step. Our goal is to improve collision avoidance and maintain route progress without retraining or modifying these upstream modules.

To this end, we introduce a lightweight \emph{safety adapter} that sits between the nominal controller and the actuator. The adapter is a Control Barrier Function (CBF) filter whose barrier gains are dynamically adjusted by a reinforcement learning (RL) policy. Perception, planning, and the nominal steering controller remain unchanged throughout training and evaluation, ensuring that any observed improvements in safety or efficiency can be directly attributed to the adaptive CBF layer.

At each decision step $t$ the deployed stack is:
\begin{equation}
\begin{aligned}
\text{Sensors} \;\rightarrow\; \text{Collaborative Perception} \;\rightarrow\; \text{Planner} \\
\rightarrow\; \text{Nominal Controller} \;\rightarrow\; \text{RL-adaptive CBF} \;\rightarrow\; \text{Actuation}.
\end{aligned}
\end{equation}
CARLA advances in synchronous mode at a fixed physics rate (e.g. 20 Hz), but the agent may skip intermediate frames and reuse the previous control command. RL transitions and CBF optimization are computed only on \emph{decision frames}, so the RL step count corresponds to decision frames rather than raw simulator ticks.

%------------------------------------------------------------------
\subsection{Nominal CoDriving Control Interface}
%------------------------------------------------------------------

Let $v_t$ denote the ego speed and $v_t^{\text{des}}$ the desired speed computed by the CoDriving controller from the planned waypoints and current target point. We convert this speed target into a nominal longitudinal acceleration:
\begin{equation}
a_t^{\text{nom}} = \frac{v_t^{\text{des}} - v_t}{\Delta t},
\end{equation}
where $\Delta t$ is the interval between control updates, not the simulator physics tick. In our experiments, CARLA runs at 20\,Hz while control is updated every 4 ticks, so we use $\Delta t=0.2$\,s in the CBF update.

Steering is generated by the nominal PID lateral controller and passed through unchanged; the safety layer modifies only longitudinal behavior. This keeps the optimization lightweight and allows us to isolate the effect of adaptive CBF gains.

%------------------------------------------------------------------
\subsection{Adaptive CBF Safety Layer}
%------------------------------------------------------------------

The CBF filter receives the nominal acceleration $a_t^{\text{nom}}$, the current ego state, and the states of nearby traffic participants, and returns a safe acceleration $a_t$ that stays as close to nominal as possible while respecting barrier-based safety constraints. We describe the barrier construction, the quadratic program, and the fallback behavior.

\subsubsection{Relative Kinematics and Barrier Definition}

For each actor $i$, we apply a range filter and then a path-aware relevance filter, focusing the CBF on actors that are likely to affect the ego vehicle's near-term trajectory. Let $p_{\text{ego}}$ and $v_{\text{ego}}$ denote the ego vehicle's position and velocity in the world frame, and $p_i$ and $v_i$ those of actor $i$. The line-of-sight unit vector from the ego to actor $i$ is:
\begin{equation}
\hat{u}_i = \frac{p_i - p_{\text{ego}}}{\|p_i - p_{\text{ego}}\| + \varepsilon},
\end{equation}
where $\varepsilon > 0$ is a small constant to prevent division by zero.
The Euclidean distance is $d_i = \|p_i - p_{\text{ego}}\|$. The signed relative velocity projected onto this line of sight is:
\begin{equation}
v_i^{\text{los}} = \hat{u}_i^\top (v_i - v_{\text{ego}}),
\end{equation}
where a negative value indicates that the actor is approaching the ego. The nonnegative closing speed is:
\begin{equation}
v_i^{\text{close}} = \max\!\left(0,\; -v_i^{\text{los}}\right).
\end{equation}

Both the ego and each actor are enclosed by a circular footprint whose radius is derived from the CARLA bounding box half-extents as $r = \sqrt{e_x^2 + e_y^2}$, where $e_x$ and $e_y$ are the longitudinal and lateral half-extents of the actor's axis-aligned bounding box. We write $r_{\text{ego}}$ for the enclosing radius of the ego vehicle and $r_i$ for that of the actor $i$. With a geometric margin $m > 0$ and a maximum deceleration parameter of braking  $a_{\text{brake}}$, the safety distance and the barrier value are:
\begin{equation}
d_i^{\text{safe}} = r_{\text{ego}} + r_i + m,
\end{equation}
\begin{equation}
h_i = d_i - d_i^{\text{safe}} - \frac{(\tilde v_i^{\text{close}})^2}{2\,a_{\text{brake}}}, \qquad
\tilde v_i^{\text{close}} = v_i^{\text{close}} + a_i^{\max}\,\Delta t.
\end{equation}

Each term in $h_i$ has a clear physical meaning. The first, $d_i$, is the current center-to-center distance. The second, $d_i^{\text{safe}}$, covers the physical sizes of both agents plus the clearance margin. The third, $(\tilde v_i^{\text{close}})^2 / (2\, a_{\text{brake}})$, is the stopping distance that ego would need if it braked at its maximum deceleration $a_{\text{brake}}$ from a closing speed that bounds the actor's own acceleration over the decision interval: rather than assuming actor velocity is exactly constant, $\tilde v_i^{\text{close}}$ inflates the raw closing speed by the actor's maximum plausible acceleration toward the ego, $a_i^{\max}$, applied over $\Delta t$. This bounds the actor-acceleration risk that a pure constant-velocity model would otherwise ignore, which matters most for pedestrians and cyclists, whose accelerations are a larger fraction of their speed than a vehicle's. When $h_i \ge 0$, there is still enough room to stop before reaching the actor $i$ even under this worst-case acceleration bound; when $h_i < 0$, the ego is already inside the braking zone and the risk of collision is high. Actors who are far away or moving apart ($v_i^{\text{close}} = 0$) produce large positive $h_i$ and do not constrain the ego.


The path-aware filter is defined in the ego frame. Actors who are sufficiently behind the ego are ignored, and actors far from the ego lane corridor are ignored unless they are likely to cross into the ego path within a short time horizon. This reduces unnecessary braking caused by traffic in adjacent lanes or in the opposite-direction, while still preserving responsiveness to relevant crossing interactions.

\subsubsection{CBF Quadratic Program}

At each decision frame, we solve a small QP over the scalar longitudinal acceleration $a_t$ and one slack variable $s_i \ge 0$ per active actor:
\begin{align}
\min_{a_t,\,\{s_i\}} \quad & \left(a_t - a_t^{\text{nom}}\right)^2 + w_s \sum_i s_i^2, \label{eq:qp-obj}\\
\text{s.t.} \quad & a_{\min} \le a_t \le a_{\max}, \label{eq:qp-acc}\\
& \alpha_i\, a_t + s_i \ge \beta_i, \quad \forall\, i. \label{eq:qp-cbf}
\end{align}

The quadratic objective~\eqref{eq:qp-obj} keeps the safe acceleration as close to the nominal command as possible, while the penalty weight $w_s$ on slack discourages but does not prohibit soft constraint violations, making the filter robust to infeasible configurations that can arise in dense traffic. Acceleration bounds~\eqref{eq:qp-acc} enforce physical actuation limits. Each CBF constraint~\eqref{eq:qp-cbf} encodes, for actor $i$, the requirement that the barrier value does not decrease faster than a rate governed by adaptive gain $\gamma_i$.

The linearized coefficients are:
\begin{align}
\alpha_i &= -\tfrac{1}{2}\,\bar\eta_i\,\Delta t^2, \\
\beta_i  &= -\gamma_i\, h_i \;-\; v_i^{\text{los}}\,\Delta t, \\
\eta_i   &= \hat{e}_{\text{ego}}^\top \hat{u}_i,
\end{align}
where $\hat{e}_{\text{ego}}$ is the unit vector along the direction of the ego vehicle and $\eta_i$ is the direction alignment, defined as the cosine of the angle between the ego's forward direction and the line of sight to actor $i$. An actor directly ahead has $\eta_i \approx 1$, making the constraint more restrictive; actors behind or to the side have smaller or negative $\eta_i$ and impose weaker longitudinal restrictions.

At a near-perpendicular crossing, $\eta_i \approx 0$, which collapses $\alpha_i$ to zero and degenerates constraint~\eqref{eq:qp-cbf} to $s_i \ge \beta_i$: the QP can satisfy it purely through slack, with no actual braking, exactly when a crossing pedestrian or cyclist most needs longitudinal protection. We guard against this by flooring the alignment magnitude used in $\alpha_i$ once the actor is already inside the danger zone:
\begin{equation}
\bar\eta_i =
\begin{cases}
\operatorname{sign}(\eta_i)\,\eta_{\text{floor}}, & h_i < 0 \text{ and } |\eta_i| < \eta_{\text{floor}}, \\
\eta_i, & \text{otherwise},
\end{cases}
\end{equation}
where $\eta_{\text{floor}} > 0$ is a small constant. The floor only activates once the barrier is already violated ($h_i < 0$); it never flips the sign of $\eta_i$ and leaves well-aligned actors ($|\eta_i| \ge \eta_{\text{floor}}$) untouched, so the constraint keeps requiring the QP to trade off against a crossing actor instead of discharging the requirement through slack alone.

\textit{Derivation of $\alpha_i$.} The coefficient $\alpha_i$ is the sensitivity of $h_i$ to the chosen ego acceleration $a_t$. This sensitivity has two channels: (i) the direct kinematic shift in relative position over $\Delta t$, and (ii) the change in closing speed induced by $a_t$, which feeds back into the braking-margin term of $h_i$. We linearize using channel (i) only, giving $\alpha_i = -\tfrac{1}{2}\bar\eta_i \Delta t^2$, and treat channel (ii) as a zero-order hold, evaluated at the current closing speed rather than expanded into the coefficient. This is a deliberate one-step approximation rather than an omission: because the filter replans every $\Delta t{=}0.2$\,s, any error from holding the closing speed fixed within one step is corrected at the next decision when $h_i$ is recomputed from the actor's realized velocity, following the same receding-horizon argument used to justify one-step linearizations in MPC-CBF filters more broadly. Section~\ref{results} reports how often this approximation is actually load-bearing, via the measured slack-activation rate.

An important detail of the implementation is that the barrier $h_i$ uses the nonnegative closing speed $v_i^{\text{close}}$ (so receding actors contribute zero braking penalty), while the QP coefficient $\beta_i$ uses the signed velocity $v_i^{\text{los}}$ (so receding actors actively relax the constraint). This asymmetry is deliberate: it prevents the barrier from becoming artificially negative for non-threatening actors while allowing the QP to account for the full relative dynamics correctly.

We solve QP with OSQP\cite{stellato2020osqp}. If the solver fails, or if no actor remains after range- and path-aware filtering, the nominal command is used directly. After solving, the safe desired speed is obtained from the QP-optimal acceleration $a_t^{*}$:
\begin{equation}
v_t^{\text{safe}} = \max\!\left(0,\; v_t + a_t^{*}\,\Delta t\right),
\end{equation}
The throttle and brake commands are recomputed from $v_t^{\text{safe}}$ via the standard V2Xverse PID speed controller.

%------------------------------------------------------------------
\subsection{Adaptive Gamma Policy}
%------------------------------------------------------------------

A fixed $\gamma$ imposes a single global safety-performance trade-off in all scenes, which is suboptimal in heterogeneous urban traffic. The core idea of our approach is to have a learned policy choose $\gamma$ for each actor class at each decision step, conditioned on the current traffic context.

\subsubsection{State and Action Spaces}

The policy observes an 11-dimensional feature vector summarizing the ego state and the nearest actor of each class:
\begin{align}
o_t = \bigl[\,&v_t,\; \Delta v_t,\; d_v,\; c_v,\; \eta_v,\; d_p,\; c_p,\; \eta_p,\; d_b,\; c_b,\; \eta_b\,\bigr],
\end{align}
where $\Delta v_t = v_t^{\text{des}} - v_t$ is the speed error, and for each actor class $k \in \{v, p, b\}$ (vehicle, pedestrian, bicycle), $(d_k, c_k, \eta_k)$ are the distance, closing speed and heading alignment of the nearest actor of that class. When no actor of a given class is present, the corresponding entries are set to sentinel values (large distance, zero closing speed, zero alignment). This compact representation captures the most critical safety information while remaining low-dimensional enough for sample-efficient on-policy learning.

The policy outputs a normalized 3-D action $\tilde{a}_t \in [-1,1]^3$, which is affinely rescaled to class-specific CBF gains:
\begin{equation}
\gamma_t^{(k)} = \frac{\tilde{a}_t^{(k)} + 1}{2}\left(\gamma_{\max}^{(k)} - \gamma_{\min}^{(k)}\right) + \gamma_{\min}^{(k)}.
\end{equation}
Limits $[\gamma_{\min}^{(k)},\, \gamma_{\max}^{(k)}]$ are hyperparameters that constrain the policy to a safe operating range.

\subsubsection{Network Architecture}

Both the actor and critic are implemented as multi-layer perceptrons with two hidden layers of 64 units each and Tanh activations. The actor outputs the mean of a diagonal Gaussian whose log standard deviation is a learnable parameter; the output is passed through a tanh squashing function so that the raw actions lie in $[-1,1]$, and the log-probability is corrected with the standard change-of-variables Jacobian. During training, actions are sampled stochastically for exploration; at evaluation time, the mean is used deterministically. The critic maps the same observation to a scalar state-value estimate. A separate reward critic is used; cost returns are computed as discounted sums directly without a dedicated cost value function.

\subsubsection{Reward Design}

The reward encourages the progress of the route. At each decision step, the reward is the nonnegative increase in route completion:
\begin{equation}
r_t^{\text{step}} = \max\!\left(0,\; \ell_t - \ell_{t-1}\right),
\end{equation}
where $\ell_t$ is the cumulative distance traveled along the route by step $t$. At route termination, an additional completion bonus is added:
\begin{equation}
r_T = r_T^{\text{step}} + \mathbb{I}_{\text{complete}}.
\end{equation}
where $\mathbb{I}_{\text{complete}}$ is 1 if the route was completed and 0 otherwise.
The nonnegative clamp prevents the policy from being rewarded for going backward (e.g., during a reversal maneuver). The completion indicator provides a sparse but meaningful signal that the route was finished successfully, encouraging the policy not to be so conservative that it stalls.

\subsubsection{Safety Cost}

Rather than folding safety into the reward (which conflates two competing objectives), we define a separate per-step cost that feeds into the constrained optimization. For each actor $i$ in the safety range, two risk measures are calculated.

\textbf{Proximity cost.} A piecewise-linear function of distance that ramps from 0 at an outer warning threshold to 1 at the collision envelope:
\begin{equation}
c_i^{\text{prox}} =
\begin{cases}
0, & d_i \ge d_{\text{warn}},\\[2pt]
\displaystyle\frac{d_{\text{warn}} - d_i}{d_{\text{warn}} - d_{\text{col}}}, & d_{\text{col}} < d_i < d_{\text{warn}},\\[6pt]
1, & d_i \le d_{\text{col}},
\end{cases}
\end{equation}
where $d_{\text{warn}}=10.0$\,m is a proximity warning distance beyond which no cost accrues (distinct from the CBF safe distance $d_i^{\text{safe}}$), and $d_{\text{col}}=2.5$\,m is the approximate collision-contact distance.


\textbf{Time-to-collision cost.} A linear ramp based on how soon contact would occur at the current closing speed:
\begin{align}
\text{TTC}_i &= \frac{d_i}{v_i^{\text{close}} + \varepsilon}, \\
c_i^{\text{ttc}} &= \max\!\left(0,\; 1 - \frac{\text{TTC}_i}{\tau_{\text{ttc}}}\right).
\end{align}
Here $\tau_{\text{ttc}}=3.0$\,s is a time-to-collision horizon; the cost is zero when $\text{TTC}_i \ge \tau_{\text{ttc}}$ (actor receding or far away) and increases toward 1 as imminent contact approaches.


The cost of decision-step aggregates across all actors by taking the worst case:
\begin{equation}
c_t^{\text{step}} = \max_i \left\{\max\!\left(c_i^{\text{prox}},\; c_i^{\text{ttc}}\right)\right\}.
\end{equation}
At route termination, a discrete collision penalty is added:
\begin{equation}
c_T = c_T^{\text{step}} + c_{\text{col}}\,\mathbb{I}_{\text{collision}},
\end{equation}
where $c_{\text{col}}$ is a configurable terminal penalty and $\mathbb{I}_{\text{collision}}$ is equal to 1 if a collision occurred during the route and 0 otherwise. The dense per-step cost provides a continuous learning signal for the policy to preemptively avoid risky states, while the terminal penalty ensures that actual collisions are heavily penalized even if they occur from states that appeared safe one step earlier.

%------------------------------------------------------------------
\subsection{Constrained Policy Optimization}
%------------------------------------------------------------------

We train the adaptive gamma policy using a PPO-style algorithm extended with a Lagrangian constraint on the cumulative safety cost. The total loss is:
\begin{equation}
\mathcal{L} =
\mathcal{L}_{\text{reward}}
+ \lambda\,\mathcal{L}_{\text{cost}}
- \alpha_H\,\mathcal{H}[\pi]
+ \alpha_V\,\mathcal{L}_{\text{value}},
\end{equation}
where $\alpha_H$ is the entropy bonus coefficient, $\alpha_V$ is the value-function loss coefficient, $\mathcal{H}[\pi]$ is the policy entropy (encouraging exploration), $\mathcal{L}_{\text{value}}$ is the loss of the reward-critic regression, and $\lambda$ is a dynamically adjusted Lagrangian multiplier that controls the weight of the safety cost relative to reward.


\subsubsection{Clipped Surrogates}

Let $\rho_t = \pi_\theta(a_t | o_t)\,/\,\pi_{\theta_{\text{old}}}(a_t | o_t)$ be the importance ratio and $\epsilon \in (0,1)$ the PPO clipping coefficient. The reward and cost policy-gradient surrogates are as follows:
\begin{align}
\mathcal{L}_{\text{reward}} &=
-\mathbb{E}\!\left[\min\!\left(\rho_t\, A_t,\;\text{clip}(\rho_t,\,1{-}\epsilon,\,1{+}\epsilon)\,A_t\right)\right], \\
\mathcal{L}_{\text{cost}} &=
\phantom{-}\mathbb{E}\!\left[\max\!\left(\rho_t\, C_t,\;\text{clip}(\rho_t,\,1{-}\epsilon,\,1{+}\epsilon)\,C_t\right)\right],
\end{align}
where $A_t$ are GAE-based reward advantages (normalized to zero mean and unit variance) and $C_t$ are discounted cumulative cost returns computed directly without a cost baseline and left unnormalized. The asymmetry in clipping direction is deliberate: the $\min$ in $\mathcal{L}_{\text{reward}}$ is standard PPO (preventing large reward-increasing updates), while the $\max$ in $\mathcal{L}_{\text{cost}}$ is conservative and takes the larger of the clipped and unclipped cost surrogate, preventing the optimizer from under-counting safety violations through ratio clipping.

\subsubsection{PID Lagrangian Update}

The multiplier $\lambda$ is updated after each rollout batch using a PID controller. Let $\bar{c}_u$ denote the mean safety cost per-step in the re-introduction $u$, and let $c_{\text{limit}}$ be the target cost threshold. The constraint error is as follows:

\begin{equation}
e_u = \bar{c}_u - c_{\text{limit}}.
\end{equation}
The multiplier is:
\begin{equation}
\lambda_u = \text{clip}\!\left(k_p\,e_u + k_i\,I_u + k_d\,D_u,\;\; 0,\;\; \lambda_{\max}\right),
\end{equation}
where $k_p$, $k_i$, and $k_d$ are proportional, integral, and derivative gains, $I_u$ is the running integral of past constraint errors, and $D_u$ is a finite-difference derivative term computed from recent batch-cost statistics.
The PID formulation provides smoother multiplier dynamics than a simple dual gradient ascent, which is particularly important in our setting where rollout cost estimates are noisy due to the stochasticity of CARLA traffic scenarios.



%------------------------------------------------------------------
\subsection{Training and Evaluation Protocol}
%------------------------------------------------------------------

Training is fully on-policy: the adaptive-gamma policy collects rollouts in CARLA, updates via the constrained PPO loss, and repeats. Collaborative perception, the CoDriving planner, and the nominal steering controller remain unchanged throughout; only the small actor-critic network is optimized. This makes training lightweight and ensures that any change in driving behavior is attributable solely to adaptive safety modulation. For evaluation, we load the trained policy checkpoint and select actions deterministically (using the actor mean without sampling noise). All other components are identical to training, which allows for fair comparison against fixed-CBF and unfiltered baselines.
